\begin{proof}[Proof of \cref{obj:lem:supplied-block-record-lower}]
\begin{enumerate}
\item The design assumptions identify the joint experimental law. Relabel the population as
\[
V=\{1,\ldots,n\},
\qquad
D=\operatorname{Law}(Z,W).
\]
The assignment marginal is a probability measure, so \(D\) is a probability measure. Independence and the two specified marginals give
\[
D
=
\operatorname{Law}(Z)\otimes\operatorname{Law}(W)
=
\left(\bigotimes_{j\in V}\operatorname{Bernoulli}(1/2)\right)
\otimes
\left(\bigotimes_{\substack{i,j\in V\\i\ne j}}
\operatorname{Bernoulli}(q)\right),
\]
by \cref{obj:ass:assignment,obj:ass:audit,obj:ass:design-independence}.
We work with this product design, first for the complete-record risk and then for the supplied-graph risk.
% lean: CausalSmith.Experimentation.ThinnedgraphAdditiveRiskFrontier.design_eq_thinnedDesign

\item For the complete-record case, define the block size, block count, covered fraction, and amplitude by
\[
k_0=d+1,
\qquad
r_0=\left\lfloor\frac{n}{k_0}\right\rfloor,
\qquad
\rho=\frac{r_0k_0}{n},
\qquad
h=\frac{1}{100\pi}
\sqrt{\min\left\{1,\frac{k_0^2}{n}\right\}}.
\]
Here \(2\le k_0\le n\), so \(r_0\ge1\) and every denominator is positive. The quantity under the square root is positive and at most one. Since \(\pi\ge2\),
\[
0<h\le\frac{1}{100\pi}\le\frac14,
\qquad
h^2=
\frac{1}{10000\pi^2}
\min\left\{1,\frac{k_0^2}{n}\right\}.
\]
To check coverage, define the integer remainder by
\[
\operatorname{rem}_0=n-r_0k_0.
\]
Euclidean division and \(r_0\ge1\) imply
\[
0\le\operatorname{rem}_0<k_0\le r_0k_0,
\qquad
n=r_0k_0+\operatorname{rem}_0\le2r_0k_0.
\]
Consequently,
\[
\frac12\le\rho\le1,
\qquad
\Delta:=\rho h>0.
\]
% lean: CausalSmith.Experimentation.ThinnedgraphAdditiveRiskFrontier.completeBlockLowerAmplitude_properties
% lean: CausalSmith.Experimentation.ThinnedgraphAdditiveRiskFrontier.completeBlock_coverage_ge_half

\item Construct two priors on fixed schedules. Define the labeled blocks and their common graph by
\[
\mathcal B_v=\{(v-1)k_0+1,\ldots,vk_0\},
\quad v=1,\ldots,r_0,
\qquad
G_{n,d}
=
\bigcup_{v=1}^{r_0}
\{(j,i):i,j\in\mathcal B_v,\ j\ne i\}.
\]
Each block has \(k_0\) members. Both off-diagonal degrees are \(k_0-1=d\) on the blocks and zero on the remaining units.

Use the baseline density of \cref{obj:lem:baseline-translation-affinity},
\[
f(w)=
\begin{cases}
4\cos^2(2\pi w),& |w|\le1/4,\\
0,& |w|>1/4.
\end{cases}
\]
It is nonnegative and integrable, and
\[
\int_{\mathbb R}f(w)\,dw
=
\int_{-1/4}^{1/4}4\cos^2(2\pi w)\,dw
=
\left[2w+\frac{\sin(4\pi w)}{2\pi}\right]_{-1/4}^{1/4}
=1.
\]
Define the baseline-vector law by
\[
U=(U_1,\ldots,U_{r_0}),
\qquad
\operatorname{Law}(U)=\bigotimes_{v=1}^{r_0}f(w)\,dw.
\]
This is a probability law, and \(|U_v|\le1/4\) for every \(v\) almost surely.

For \(\sigma\in\{-1,1\}\) and \(U\in\mathbb R^{r_0}\), let
\(\theta^{\mathrm{cb}}_{\sigma,h}(U)\) be the schedule with graph \(G_{n,d}\) and coefficients
\[
a_i=
\begin{cases}
U_v-\sigma h/2,&i\in\mathcal B_v,\\
0,&i>r_0k_0,
\end{cases}
\qquad
t_i=
\begin{cases}
\sigma h/k_0,&i\le r_0k_0,\\
0,&i>r_0k_0,
\end{cases}
\qquad
b_{ij}=\frac{\sigma h}{k_0}
\quad\bigl((j,i)\in G_{n,d}\bigr).
\]
For a supported baseline draw and \(i\in\mathcal B_v\),
\[
|a_i|+|t_i|+\sum_{j:(j,i)\in G_{n,d}}|b_{ij}|
=
|U_v-\sigma h/2|+\frac{h}{k_0}
+d\frac{h}{k_0}
\le\frac14+\frac h2+h
\le\frac58\le1.
\]
The corresponding expression is zero for a padding unit. Thus these schedules belong almost surely to the class in \cref{obj:def:model}.

For a binary assignment \(z\in\{0,1\}^{V}\), define
\[
X_j(z)=2z_j-1,
\qquad j\in V.
\]
The additive outcome formula gives
\[
Y_i(z)=
\begin{cases}
\displaystyle
U_v-\frac{\sigma h}{2}
+\frac{\sigma h}{k_0}\sum_{j\in\mathcal B_v}z_j
=
U_v+\frac{\sigma h}{2k_0}
\sum_{j\in\mathcal B_v}X_j(z),
&i\in\mathcal B_v,\\[6pt]
0,&i>r_0k_0.
\end{cases}
\]
In particular,
\[
Y_i(\mathbf1)-Y_i(\mathbf0)=
\begin{cases}
\sigma h,&i\le r_0k_0,\\
0,&i>r_0k_0,
\end{cases}
\qquad
\tau\bigl(\theta^{\mathrm{cb}}_{\sigma,h}(U)\bigr)
=
\sigma\frac{r_0k_0h}{n}
=\sigma\Delta.
\]
Define the priors by
\[
\Pi^{\mathrm{cb}}_{\sigma,h}
=
\operatorname{Law}\bigl(\theta^{\mathrm{cb}}_{\sigma,h}(U)\bigr).
\]
Under each prior the schedule is drawn before, and independently of, the experiment. The recording map is jointly measurable: its graph and assignment coordinates are finite, and its outcome coordinates are affine functions of \(U\) for each design draw.
% lean: CausalSmith.Experimentation.ThinnedgraphAdditiveRiskFrontier.completeBlock_support
% lean: CausalSmith.Experimentation.ThinnedgraphAdditiveRiskFrontier.completeBlock_record_measurable

\item Define the complete-record mixture laws by
\[
\mathsf Q_+
=
\int P_{\theta,q}\,\Pi^{\mathrm{cb}}_{1,h}(d\theta),
\qquad
\mathsf Q_-
=
\int P_{\theta,q}\,\Pi^{\mathrm{cb}}_{-1,h}(d\theta).
\]
The support, constant-target, and measurability requirements of
\cref{obj:lem:full-record-two-prior-risk} have just been verified. It therefore supplies
\[
R_{n,d}(q)
\ge
\frac{\Delta^2}{2}
\left(1-\operatorname{TV}(\mathsf Q_+,\mathsf Q_-)\right).
\]
% lean: CausalSmith.Experimentation.ThinnedgraphAdditiveRiskFrontier.full_record_two_prior_risk

\item To bound this testing distance, retain the entire design draw \((z,w)\) and the vector \(y\in\mathbb R^{r_0}\) of distinct block responses. Define the shifts and response densities by
\[
\delta_{\sigma,v}(z)
=
\frac{\sigma h}{2k_0}
\sum_{j\in\mathcal B_v}X_j(z),
\qquad
\sigma\in\{-1,1\},\quad v=1,\ldots,r_0,
\]
\[
\lambda^{\mathrm{cb}}_{\sigma,h}(y\mid z,w)
=
\prod_{v=1}^{r_0}f\bigl(y_v-\delta_{\sigma,v}(z)\bigr).
\]
For every design draw these densities are nonnegative and normalized:
\[
\int_{\mathbb R^{r_0}}
\lambda^{\mathrm{cb}}_{\sigma,h}(y\mid z,w)\,dy
=
\prod_{v=1}^{r_0}
\int_{\mathbb R}f\bigl(y_v-\delta_{\sigma,v}(z)\bigr)\,dy_v
=1.
\]
Indeed, translating \(U_v\) by \(\delta_{\sigma,v}(z)\) gives the density
\(f(y_v-\delta_{\sigma,v}(z))\); taking the product gives the displayed vector density.

Define the joint laws by
\[
\mathsf J_\pm(dz,dw,dy)
=
D(dz,dw)\,
\lambda^{\mathrm{cb}}_{\pm1,h}(y\mid z,w)\,dy.
\]
These are precisely the laws of the full design draw together with
\((U_v+\delta_{\pm1,v}(Z))_{v=1}^{r_0}\).

For an audit array \(w\), define the retained graph by
\[
H(w)_{ij}
=
\mathbf1\{(j,i)\in G_{n,d}\}\,w_{ij},
\qquad i\ne j.
\]
The common recording map is
\[
\mathcal C_{\mathrm{rec}}((z,w),y)
=
\left(
H(w),\,z,\,
\left(
\begin{cases}
y_v,&i\in\mathcal B_v,\\
0,&i>r_0k_0
\end{cases}
\right)_{i\in V}
\right).
\]
The outcome identity in the preceding construction proves
\[
O
=
\mathcal C_{\mathrm{rec}}
\left((Z,W),(U_v+\delta_{\sigma,v}(Z))_{v=1}^{r_0}\right),
\qquad
\mathsf Q_\pm=(\mathcal C_{\mathrm{rec}})_\#\mathsf J_\pm.
\]
Here the subscript \(\#\) denotes the pushforward of a measure by the displayed map. The map is measurable, and its pushforward preserves the labeled retained graph and every treatment coordinate. Taking preimages of measurable events yields
\[
\operatorname{TV}(\mathsf Q_+,\mathsf Q_-)
\le
\operatorname{TV}(\mathsf J_+,\mathsf J_-).
\]
% lean: CausalSmith.Experimentation.ThinnedgraphAdditiveRiskFrontier.completeBlock_mixture_density_representation
% lean: CausalSmith.Experimentation.ThinnedgraphAdditiveRiskFrontier.completeBlock_mixture_tv_population_le

\item Define the conditional squared Hellinger energy by
\[
\mathcal E_{\mathrm{cb}}(z,w)
=
\int_{\mathbb R^{r_0}}
\left(
\sqrt{\lambda^{\mathrm{cb}}_{1,h}(y\mid z,w)}
-
\sqrt{\lambda^{\mathrm{cb}}_{-1,h}(y\mid z,w)}
\right)^2dy.
\]
The product-translation bound in
\cref{obj:lem:baseline-translation-affinity} gives
\[
\mathcal E_{\mathrm{cb}}(z,w)
\le
\sum_{v=1}^{r_0}
4\pi^2
\bigl(\delta_{1,v}(z)-\delta_{-1,v}(z)\bigr)^2.
\]

Write \(X_j=X_j(Z)\) for the random treatment signs. Under the assignment marginal of \(D\),
\[
\mathbb E_D[X_j]=0,
\qquad
\mathbb E_D[X_j^2]=1,
\qquad
\mathbb E_D[X_iX_j]=0\quad(i\ne j).
\]
Expanding the square therefore gives, for every block,
\[
\mathbb E_D
\left[\left(\sum_{j\in\mathcal B_v}X_j\right)^2\right]
=
\sum_{j\in\mathcal B_v}\mathbb E_D[X_j^2]
+
\sum_{\substack{i,j\in\mathcal B_v\\i\ne j}}
\mathbb E_D[X_iX_j]
=k_0.
\]
Since
\[
\delta_{1,v}(Z)-\delta_{-1,v}(Z)
=
\frac{h}{k_0}\sum_{j\in\mathcal B_v}X_j,
\]
it follows that
\[
\mathbb E_D
\left[
\bigl(\delta_{1,v}(Z)-\delta_{-1,v}(Z)\bigr)^2
\right]
=
\frac{h^2}{k_0}.
\]
Define the averaged energy by
\[
\mathcal E_{\mathrm{joint}}
=
\int\mathcal E_{\mathrm{cb}}(z,w)\,D(dz,dw).
\]
Integrating the conditional bound and using \(r_0k_0\le n\) gives
\[
\mathcal E_{\mathrm{joint}}
\le
4\pi^2h^2\frac{r_0}{k_0}
\le
4\pi^2h^2\frac{n}{k_0^2}.
\]

The common-marginal identity in
\cref{obj:lem:baseline-translation-affinity} identifies
\(\mathcal E_{\mathrm{joint}}\) as the unhalved squared Hellinger distance between \(\mathsf J_+\) and \(\mathsf J_-\).
For completeness, the sharper total-variation comparison used here follows directly from Cauchy--Schwarz. Define the common reference measure by
\[
\mu_{\mathrm{cb}}(dz,dw,dy)=D(dz,dw)\,dy.
\]
Normalization and the definition of the energy imply
\[
\int
\left(
\sqrt{\lambda^{\mathrm{cb}}_{1,h}}
+
\sqrt{\lambda^{\mathrm{cb}}_{-1,h}}
\right)^2d\mu_{\mathrm{cb}}
=
4-\mathcal E_{\mathrm{joint}}.
\]
Consequently,
\[
\begin{aligned}
\operatorname{TV}(\mathsf J_+,\mathsf J_-)
&=
\frac12\int
\left|\lambda^{\mathrm{cb}}_{1,h}
-\lambda^{\mathrm{cb}}_{-1,h}\right|d\mu_{\mathrm{cb}}\\
&\le
\frac12
\sqrt{\mathcal E_{\mathrm{joint}}}
\sqrt{4-\mathcal E_{\mathrm{joint}}}\\
&=
\sqrt{\mathcal E_{\mathrm{joint}}
\left(1-\mathcal E_{\mathrm{joint}}/4\right)}
\le
\sqrt{\mathcal E_{\mathrm{joint}}}.
\end{aligned}
\]
The densities in these integrals are evaluated at \((y\mid z,w)\). Together with the recording-map bound, this proves
\[
\operatorname{TV}(\mathsf Q_+,\mathsf Q_-)
\le
\sqrt{4\pi^2h^2\frac{n}{k_0^2}}.
\]
% lean: CausalSmith.Experimentation.ThinnedgraphAdditiveRiskFrontier.completeBlock_shift_separation_moment
% lean: CausalSmith.Experimentation.ThinnedgraphAdditiveRiskFrontier.completeBlock_response_energy_population_le
% lean: CausalSmith.Experimentation.ThinnedgraphAdditiveRiskFrontier.completeBlock_densityJoint_tv_le
% lean: CausalSmith.Experimentation.ThinnedgraphAdditiveRiskFrontier.completeBlock_mixture_tv_population_le

\item The amplitude calibration now gives
\[
4\pi^2h^2\frac{n}{k_0^2}
=
\min\left\{1,\frac{k_0^2}{n}\right\}
\frac{n}{2500k_0^2}
\le\frac1{2500},
\]
and hence
\[
\operatorname{TV}(\mathsf Q_+,\mathsf Q_-)\le\frac1{50}.
\]
Coverage and positivity give the separation comparisons
\[
\frac h2\le\rho h=\Delta,
\qquad
\frac{h^2}{4}\le\Delta^2.
\]
The exact squared-amplitude identity also gives
\[
\frac{1}{160000\pi^2}
\min\left\{1,\frac{(d+1)^2}{n}\right\}
=\frac{h^2}{16}.
\]
Substituting these comparisons into the complete-record two-prior bound yields
\[
\begin{aligned}
R_{n,d}(q)
&\ge
\frac{\Delta^2}{2}
\left(1-\operatorname{TV}(\mathsf Q_+,\mathsf Q_-)\right)\\
&\ge
\frac{h^2}{8}\left(1-\frac1{50}\right)
\ge\frac{h^2}{16}\\
&=
\frac{1}{160000\pi^2}
\min\left\{1,\frac{(d+1)^2}{n}\right\}.
\end{aligned}
\]
% lean: CausalSmith.Experimentation.ThinnedgraphAdditiveRiskFrontier.completeBlock_testing_distance
% lean: CausalSmith.Experimentation.ThinnedgraphAdditiveRiskFrontier.completeBlock_risk_lower_assembly

\item For the supplied-graph case, use the same amplitude and priors. Their support and target properties are
\[
\Pi^{\mathrm{cb}}_{\sigma,h}(\mathcal M_{n,d})=1,
\qquad
\tau\bigl(\theta^{\mathrm{cb}}_{\sigma,h}(U)\bigr)=\sigma\Delta,
\qquad
G\bigl(\theta^{\mathrm{cb}}_{\sigma,h}(U)\bigr)=G_{n,d}.
\]
Fix any measurable randomized rule \(T^G\) admitted by
\cref{obj:def:supplied-risk}. Define its rule with the graph fixed, and its supplied-graph worst-case risk, by
\[
T_{\mathrm{fix}}(o,\xi)=T^G(o,G_{n,d},\xi),
\qquad \xi\in[0,1],
\]
\[
\mathcal W^G(T^G)
=
\sup_{\theta\in\mathcal M_{n,d}}
\mathbb E_{P_{\theta,q}\otimes\nu_\xi}
\left[
\bigl(T^G(O,G(\theta),\xi)-\tau(\theta)\bigr)^2
\right],
\qquad
\nu_\xi=\operatorname{Uniform}[0,1].
\]
Inserting the constant finite graph is measurable. For every schedule in either prior,
\[
\bigl(T_{\mathrm{fix}}(O,\xi)-\tau(\theta)\bigr)^2
=
\bigl(T^G(O,G(\theta),\xi)-\tau(\theta)\bigr)^2.
\]
Integrating over the supported schedules therefore gives
\[
\int
\bigl(T_{\mathrm{fix}}(o,\xi)-\Delta\bigr)^2
\,d(\mathsf Q_+\otimes\nu_\xi)(o,\xi)
\le\mathcal W^G(T^G),
\]
\[
\int
\bigl(T_{\mathrm{fix}}(o,\xi)+\Delta\bigr)^2
\,d(\mathsf Q_-\otimes\nu_\xi)(o,\xi)
\le\mathcal W^G(T^G).
\]
These are nonnegative integrals, so the identities and inequalities also hold when the risk is infinite.
% lean: CausalSmith.Experimentation.ThinnedgraphAdditiveRiskFrontier.hardcodeSuppliedEstimator_sqLoss
% lean: CausalSmith.Experimentation.ThinnedgraphAdditiveRiskFrontier.mixture_seeded_risk_le_supplied_worst

\item Adding the common independent seed preserves total variation:
\[
\operatorname{TV}
(\mathsf Q_+\otimes\nu_\xi,\mathsf Q_-\otimes\nu_\xi)
=
\operatorname{TV}(\mathsf Q_+,\mathsf Q_-).
\]
Define the measurable testing event by
\[
\mathcal E_{\mathrm{err}}
=
\{(o,\xi):T_{\mathrm{fix}}(o,\xi)<0\}.
\]
The definition of total variation gives
\[
(\mathsf Q_+\otimes\nu_\xi)(\mathcal E_{\mathrm{err}})
+
(\mathsf Q_-\otimes\nu_\xi)(\mathcal E_{\mathrm{err}}^c)
\ge
1-\operatorname{TV}(\mathsf Q_+,\mathsf Q_-).
\]
Because \(\Delta>0\), the pointwise loss comparisons are
\[
\bigl(T_{\mathrm{fix}}(o,\xi)-\Delta\bigr)^2
\ge\Delta^2\mathbf1_{\mathcal E_{\mathrm{err}}}(o,\xi),
\]
\[
\bigl(T_{\mathrm{fix}}(o,\xi)+\Delta\bigr)^2
\ge\Delta^2\mathbf1_{\mathcal E_{\mathrm{err}}^c}(o,\xi).
\]
Integrating and using the two worst-case-risk comparisons proves
\[
2\mathcal W^G(T^G)
\ge
\Delta^2\left(1-\operatorname{TV}(\mathsf Q_+,\mathsf Q_-)\right).
\]
Division by the positive finite constant \(2\) is valid in the extended nonnegative reals. Infimizing over \(T^G\) thus yields
\[
R^G_{n,d}(q)
\ge
\frac{\Delta^2}{2}
\left(1-\operatorname{TV}(\mathsf Q_+,\mathsf Q_-)\right).
\]
The same testing-distance and separation bounds give
\[
\begin{aligned}
R^G_{n,d}(q)
&\ge
\frac{h^2}{8}\left(1-\frac1{50}\right)
\ge\frac{h^2}{16}\\
&=
\frac{1}{160000\pi^2}
\min\left\{1,\frac{(d+1)^2}{n}\right\}.
\end{aligned}
\]
The construction has at least one full block throughout the stated parameter range, and the density comparison integrates over the full design law. These arguments apply for every \(q\in[0,1]\), including both endpoints, and for every measurable randomized estimator in the two risk definitions.
% lean: CausalSmith.Experimentation.ThinnedgraphAdditiveRiskFrontier.two_prior_squared_testing
% lean: CausalSmith.Experimentation.ThinnedgraphAdditiveRiskFrontier.completeBlock_supplied_two_prior
% lean: CausalSmith.Experimentation.ThinnedgraphAdditiveRiskFrontier.completeBlock_supplied_risk_lower_assembly
\end{enumerate}
\end{proof}
